% ============================================================================
%  Section 7: Kết quả thực nghiệm và đánh giá
% ============================================================================
\section{Kết quả thực nghiệm và đánh giá}
\label{sec:results}

Chương này trình bày kết quả được tính lại trực tiếp từ ba tệp prediction trên 234 mẫu test. Các bảng metrics và latency được sinh tự động bởi script \path{scripts/generate_report_artifacts.py}; vì vậy số liệu trong báo cáo và artifact JSON có cùng nguồn.

% ============================================================================
\subsection{Kết quả trên toàn bộ mẫu và trên trường có mặt}
\label{subsec:main_results}

% Generated by scripts/generate_report_artifacts.py.
\begin{table}[htbp]
  \centering
  \caption{Kết quả trên toàn bộ 234 mẫu kiểm thử. Macro là trung bình không trọng số của bốn trường.}
  \label{tab:all-results}
  \small
  \begin{tabularx}{\textwidth}{l C{1.55cm} C{1.55cm} C{1.55cm} C{1.55cm} C{1.7cm}}
    \toprule
    \textbf{Phương pháp} & \textbf{Tên CH} & \textbf{Ngày} & \textbf{Tổng} & \textbf{Địa chỉ} & \textbf{Macro} \\
    \midrule
    \multicolumn{6}{l}{\textbf{Exact Match (EM \%) $\uparrow$}} \\
Baseline (Regex) & 25,21 & 74,79 & 54,27 & 4,70 & 39,74 \\
\model{Donut} & 3,85 & 7,69 & 7,26 & 4,70 & 5,87 \\
\model{LayoutXLM} & 29,06 & 60,68 & 69,66 & 14,53 & 43,48 \\
    \midrule
    \multicolumn{6}{l}{\textbf{Normalized Edit Similarity (NES \%) $\uparrow$}} \\
Baseline (Regex) & 40,84 & 77,01 & 67,86 & 30,29 & 54,00 \\
\model{Donut} & 8,11 & 7,69 & 7,26 & 4,70 & 6,94 \\
\model{LayoutXLM} & 56,65 & 61,75 & 74,30 & 54,73 & 61,86 \\
    \midrule
    \multicolumn{6}{l}{\textbf{Character Error Rate (CER \%) $\downarrow$}} \\
Baseline (Regex) & 91,35 & 22,99 & 37,00 & 79,49 & 57,71 \\
\model{Donut} & 140,57 & 92,31 & 92,74 & 95,30 & 105,23 \\
\model{LayoutXLM} & 45,79 & 38,25 & 25,79 & 48,74 & 39,64 \\
    \bottomrule
  \end{tabularx}
\end{table}

\begin{table}[htbp]
  \centering
  \caption{Kết quả chỉ trên các mẫu có nhãn ground-truth khác rỗng của từng trường (present-only).}
  \label{tab:present-results}
  \small
  \begin{tabularx}{\textwidth}{l C{1.55cm} C{1.55cm} C{1.55cm} C{1.55cm} C{1.7cm}}
    \toprule
    \textbf{Phương pháp} & \textbf{Tên CH} & \textbf{Ngày} & \textbf{Tổng} & \textbf{Địa chỉ} & \textbf{Macro} \\
    \midrule
    \multicolumn{6}{l}{\textbf{Exact Match (EM \%) $\uparrow$}} \\
Baseline (Regex) & 26,34 & 75,46 & 58,06 & 1,79 & 40,41 \\
\model{Donut} & 0,00 & 0,00 & 0,00 & 0,00 & 0,00 \\
\model{LayoutXLM} & 28,12 & 59,72 & 70,51 & 12,56 & 42,73 \\
    \midrule
    \multicolumn{6}{l}{\textbf{Normalized Edit Similarity (NES \%) $\uparrow$}} \\
Baseline (Regex) & 42,66 & 77,87 & 72,72 & 28,64 & 55,47 \\
\model{Donut} & 4,46 & 0,00 & 0,00 & 0,00 & 1,12 \\
\model{LayoutXLM} & 56,95 & 60,88 & 75,52 & 54,74 & 62,02 \\
    \midrule
    \multicolumn{6}{l}{\textbf{Character Error Rate (CER \%) $\downarrow$}} \\
Baseline (Regex) & 90,97 & 22,13 & 32,53 & 81,62 & 56,81 \\
\model{Donut} & 146,40 & 100,00 & 100,00 & 100,00 & 111,60 \\
\model{LayoutXLM} & 45,60 & 39,12 & 24,58 & 48,90 & 39,55 \\
    \bottomrule
  \end{tabularx}
\end{table}


\Cref{tab:all-results} phản ánh hành vi của hệ thống trên toàn bộ test set, bao gồm trường hợp trường thông tin không có nhãn. \Cref{tab:present-results} loại các ground-truth rỗng theo từng trường và được dùng làm cơ sở chính để đánh giá khả năng trích xuất.

Trên chế độ present-only, LayoutXLM đạt Macro EM 42,73\% và Macro NES 62,02\%, cao hơn Baseline tương ứng 40,41\% và 55,47\%. Khoảng cách Macro EM chỉ là 2,32 điểm phần trăm. Kết quả cho thấy checkpoint LayoutXLM hiện có tốt hơn về tổng thể, nhưng không đủ để kết luận ưu thế tuyệt đối vì hai mô hình học sâu không dùng cùng lượng và thành phần dữ liệu huấn luyện.

Donut có Macro EM present-only bằng 0\%. Mô hình không tạo giá trị chuẩn hóa khác rỗng cho \field{date}, \field{total} và \field{address}; 73/234 mẫu có dự đoán \field{store\_name} khác rỗng nhưng không mẫu có nhãn \field{store\_name} nào khớp hoàn toàn. Các điểm EM all-samples 7,69\%, 7,26\% và 4,70\% của ba trường còn lại lần lượt trùng với tỷ lệ ground-truth rỗng, nên không thể diễn giải là mô hình đã trích xuất đúng các trường đó.

\begin{figure}[htbp]
  \centering
  \includegraphics[width=0.9\textwidth]{em_present_only_comparison.png}
  \caption{So sánh Exact Match trên các mẫu có ground-truth khác rỗng. Hình được sinh từ metrics JSON.}
  \label{fig:present-only-em}
\end{figure}

% ============================================================================
\subsection{Phân tích theo trường}
\label{subsec:model_analysis}

\begin{itemize}
  \item \field{store\_name}: LayoutXLM đạt EM present-only 28,12\%, cao hơn Baseline 26,34\%. NES của LayoutXLM đạt 56,95\%, cho thấy nhiều prediction gần đúng nhưng chưa khớp hoàn toàn.
  \item \field{date}: Baseline đạt EM present-only 75,46\%, cao hơn LayoutXLM 59,72\%. Cấu trúc ngày phù hợp với regex, trong khi token classifier có thể chọn nhầm một chuỗi ngày khác trong tài liệu.
  \item \field{total}: LayoutXLM đạt EM present-only 70,51\%, so với 58,06\% của Baseline. Đây là trường LayoutXLM tạo khác biệt rõ nhất trong checkpoint hiện tại.
  \item \field{address}: Cả hai phương pháp còn yếu về exact match; LayoutXLM đạt 12,56\%, Baseline đạt 1,79\%. Tuy nhiên NES của LayoutXLM đạt 54,74\%, phản ánh nhiều trường hợp lấy được một phần địa chỉ.
\end{itemize}

Các nhận định trên là mô tả kết quả, không phải chứng minh cơ chế nhân quả. Ví dụ, chênh lệch ở \field{total} có thể liên quan đến bố cục, nhưng báo cáo chưa có ablation tách riêng text, layout và visual features.

% ============================================================================
\subsection{Latency của các thành phần đã ghi nhận}
\label{subsec:latency_evaluation}

% Generated by scripts/generate_report_artifacts.py.
\begin{table}[htbp]
  \centering
  \caption{Latency trung bình của các thành phần hợp lệ đã ghi nhận. Hai phép đo đều bắt đầu sau OCR cache và không đại diện cho latency end-to-end từ ảnh thô.}
  \label{tab:latency-comparison}
  \begin{tabular}{llr}
    \toprule
    \textbf{Phương pháp} & \textbf{Phạm vi phép đo} & \textbf{Trung bình (ms/ảnh)} \\
    \midrule
Baseline & Hậu xử lý từ OCR cache & 1,08 \\
VietOCR + \model{LayoutXLM} & Suy luận và hậu xử lý từ OCR cache & 87,11 \\
\model{Donut} & Suy luận trực tiếp từ ảnh đầu vào & 5411,17 \\
    \bottomrule
  \end{tabular}
\end{table}


\Cref{tab:latency-comparison} không phải phép so sánh end-to-end đồng nhất. Baseline và LayoutXLM bắt đầu từ OCR cache, trong khi Donut bắt đầu từ ảnh. Có thể kết luận rằng phần luật của Baseline rất nhẹ và bước suy luận LayoutXLM trong artifact nhanh hơn giải mã tự hồi quy Donut. Không thể suy ra latency sản xuất từ ảnh thô của hai pipeline OCR-based khi thiếu phép đo PaddleOCR + VietOCR cùng điều kiện phần cứng.

% ============================================================================
\subsection{Phân tích lỗi có thể tái lập}
\label{subsec:error_analysis}

Phân tích lỗi tự động sử dụng taxonomy trong \cref{tab:error-taxonomy-repro}. Mỗi prediction sai được gán đúng một nhóm theo thứ tự quy tắc trong mã nguồn. Đây là công cụ chẩn đoán nhất quán, không phải nhãn nguyên nhân hoàn hảo.

\begin{table}[htbp]
  \centering
  \caption{Taxonomy lỗi được sinh tự động từ prediction artifact.}
  \label{tab:error-taxonomy-repro}
  \begin{tabular}{lp{10cm}}
    \toprule
    \textbf{Loại lỗi} & \textbf{Tiêu chí} \\
    \midrule
    \code{EMPTY\_PRED} & Ground-truth có giá trị nhưng prediction rỗng. \\
    \code{FORMAT\_ERROR} & Prediction ngày hoặc tổng tiền không đúng schema chuẩn hóa. \\
    \code{OCR\_MISS} & Với Baseline/LayoutXLM, chuỗi gold không xuất hiện và không có cửa sổ OCR đủ gần. \\
    \code{OCR\_WRONG} & Với Baseline/LayoutXLM, OCR có cửa sổ gần gold nhưng sai ký tự. \\
    \code{POSTPROCESS\_BAD} & Prediction thô gần gold nhưng bị thay đổi sai sau chuẩn hóa. \\
    \code{MODEL\_BAD} & Các trường hợp sai còn lại của luật hoặc mô hình. \\
    \code{LABEL\_BAD} & Ground-truth rỗng nhưng hệ thống trả về một giá trị. \\
    \bottomrule
  \end{tabular}
\end{table}

Kết quả CSV cho thấy lỗi nổi bật của Donut là prediction rỗng: 216 mẫu ở \field{date}, 217 ở \field{total} và 223 ở \field{address}. Donut không được gán lỗi OCR vì pipeline không sử dụng OCR độc lập. Với LayoutXLM, \code{EMPTY\_PRED}, \code{OCR\_MISS} và \code{MODEL\_BAD} đều xuất hiện; điều này phù hợp với nhận định rằng pipeline phụ thuộc vào OCR, nhưng taxonomy tự động chưa đủ để định lượng chính xác phần hao hụt do OCR.

Các khái niệm như hallucination hoặc lỗi BIO phân mảnh chỉ được dùng trong phân tích ví dụ định tính nếu có đối chiếu ảnh, chuỗi sinh và token labels. Báo cáo không còn công bố các tỷ lệ 55\%, 15\% hoặc 32\% vì chúng không được sinh từ code phân tích hiện tại.

% ============================================================================
\subsection{Oracle OCR và câu hỏi về ảnh hưởng của OCR}
\label{subsec:oracle_ocr_proposal}

RQ3 yêu cầu đo mức ảnh hưởng của OCR lên LayoutXLM. Thực nghiệm cần thiết là huấn luyện và đánh giá cùng pipeline với văn bản/hộp bao ground-truth, sau đó so sánh với OCR thực tế. Nhánh \code{oracle\_ocr} đã được dự kiến trong dataset code nhưng chưa có artifact kết quả. Vì vậy báo cáo chỉ ghi nhận sự phụ thuộc định tính và để câu hỏi định lượng này cho nghiên cứu tiếp theo.

% ============================================================================
\subsection{Demo ứng dụng}
\label{subsec:demo}

Repository có ứng dụng Gradio cho phép chọn phương pháp, tải ảnh và xem JSON cùng latency của lượt chạy. Ứng dụng được khởi chạy bằng:
\begin{lstlisting}[language=bash]
python -m receipt_ie.app.gradio_app
\end{lstlisting}
Demo hỗ trợ kiểm tra trực quan, nhưng không thay thế protocol đánh giá trên test set.
